\subsection*{元数学}

无矛盾性一直被视为一切数学的必要条件，且自亚里士多德时代起，逻辑学已发展到足以认识到：从矛盾的理论中可以推出任何命题。自古以来就被视为不可或缺的“存在性”证明，显然仅仅是为了保证引入新概念不会导致矛盾，尤其是当该概念过于复杂，无法直接归于“直觉”之下时。我们已经看到，随着十九世纪公理化观点的出现，这种对无矛盾性证明的要求变得愈发迫切，以及算术“模型”的构造如何回应了这一要求。但算术本身难道不可能矛盾吗？在十九世纪末之前，这样的问题是不可能被提出的；整数在人们看来是如此地属于直觉中最确定的部分。然而在“悖论”出现之后，一切似乎都变得可疑，它们所造成的不安全感，可以理解地促使数学家们在大约1900年时更仔细地审视算术的一致性问题，以期至少将经典数学从沉船中拯救出来。因此，这一问题正是希尔伯特在1900年国际数学家大会著名演讲中所列出的第二个问题。{[}31{]}，第229-301页）。在此过程中，他提出了一条将会产生深远影响的新原则。在传统逻辑中，概念的无矛盾性仅仅使该概念成为“可能的”，而对希尔伯特而言，这等同于概念的存在（至少在涉及公理化定义的数学概念时是如此）。这一观点显然意味着，必须先验地证明一个数学理论的一致性，然后才能合法地发展该理论；这无疑正是庞加莱（H. Poincaré）对该原则的理解，他抓住希尔伯特的这一想法，将其作为抨击形式主义者的便利工具，并带着恶意的满足感指出，形式主义者在那个时期距离满足这一条件还有多远（{[}33 c{]}, 第163页）。我们将在后文中看到希尔伯特是如何接受这一挑战的。但首先应当指出，在他的影响以及庞加莱的影响下，后者提出的要求在很长一段时间内被毫无保留地接受，无论是形式主义者还是他们的反对者都是如此。一个后果是，即使在形式主义者中间也广泛存在一种信念，即希尔伯特的证明论是数学的一个组成部分，并且构成了数学本身不可或缺的绪论。我们在引言中已经解释了为什么这一教条在我们看来是没有根据的（*），并且我们认为，元数学在逻辑和数学论述中的作用可以而且应当仅限于处理缩写符号和演绎准则的那一非常基础的部分。与庞加莱所断言的不同，这里的问题并非“声称免于矛盾”，而是像阿达马那样认为，一致性即使无法证明，也是可以确定的。{[}35{]}，第270页）。

我们还需简要概述希尔伯特及其学派的努力历程。尽管本系列未涉及证明论，但追溯思想演变的轮廓并非毫无意义——这些思想最终导向了哥德尔的否定性结果，并事后证明了

阿达马的怀疑态度，同时也包括所有那些源于对数学推理机制理解所取得的进展，这些进展已将现代元数学提升为一门具有不可否认重要性的自主科学。

1904年，在国际大会的一次演讲中（{[}30{]}（第247-261页），希尔伯特着手处理算术一致性的问题。首先，他确立了一致性证明不能通过诉诸“模型”（*）来获得，并大致勾勒了另一种方法的原理。他提议将形式化算术中的真命题视为无意义的符号组合，并表明，通过运用支配这些组合的形成与并置的规则，永远无法得到一个既是真命题而其否定也是真命题的组合。他为一种比算术形式体系更简单的形式体系勾勒了这种性质的证明；但正如H.庞加莱不久后所观察到的（{[}33 c{]}，第185页），这个证明本质上使用了归纳原理，因此似乎建立在恶性循环之上。希尔伯特没有立即回应这一批评，十五年过去了，才有人尝试发展他的思想。直到1917年，（受回答直觉主义者攻击的愿望所驱使）他才再次投身于数学基础问题，从那时起这一问题一直占据着他的精力，直到他科学生涯的结束。在他关于这一主题的工作中，时间跨度从1920年到大约1930年，并且有一整批年轻数学家（阿克曼、贝尔奈斯、埃尔布朗、冯·诺依曼）积极参与其中，希尔伯特逐渐以更精确的方式提炼出他的“证明论”原理。他含蓄地认识到庞加莱批评的合理性，承认在元数学中，所使用的算术论证只能基于我们对整数的直觉（而非形式化的算术）。因此，将这些论证限制在直觉主义者所允许的“有限程序”（finite Prozesse）类型之内就显得至关重要。例如，反证法不能确立某个组合或组合序列的元数学存在性；必须给出一个明确的构造法则（†）。此外，希尔伯特在两个方向上扩大了他最初的纲领；他不仅着手解决算术的一致性问题，还希望证明实数理论甚至集合理论的一致性（*）。在一致性问题之外，又增加了公理独立性、范畴性和判定问题。我们将对这些不同的问题作简要回顾，并提及它们所引发的主要研究。

一个命题系统 \(A_{1}, A_{2}, \ldots, A_{n}\) 的独立性证明在于表明：对于每个指标 i， \(A_{i}\) 在理论 \(G_{i}\) 中不是定理，该理论是将 \(A_{j}\) （ \(j \neq i\) ）取作公理而得到的。我们只需要找到一个不矛盾的理论 \(G_{i}^{\prime}\) ，在其中 \(A_{j} (j \neq i)\) 是定理，连同“非 \(A_{i}\) ”；因此，我们可以根据是否假定某些理论（如算术或集合论）是一致的，从两个方面来考虑这个问题。在后一种情况下，我们面临的是“绝对”一致性的问题。另一方面，第一类问题，如同“相对”一致性问题一样，需要通过构造适当的“模型”来解决，而且许多这类性质的证明早在数学完全形式化之前就已经被设计出来了。我们可以举出非欧几里得几何的模型，以及希尔伯特在《几何基础》中处理的初等几何公理独立性问题作为例子。 {[}30{]}斯坦尼茨关于代数公理化的研究，以及豪斯多夫及其后继者关于拓扑公理化的研究。

一个理论 \(\mathfrak{C}\) 被称为范畴的，如果对于 \(\mathfrak{C}\) 中每一个除了 \(\mathfrak{C}\) 的常量外不包含任何其他字母的命题 A，两个命题 A 与（非 A）中有一个是 \(\mathfrak{C}\) 中的定理 ( \(\dagger\) )。除了一些极其初步的、其范畴性易于证明的形式系统（{[}32{]}，第35页；参见第一章附录练习7），该领域获得的结果本质上是否定的。其中最重要的是由K.哥德尔所证明的，他指出如果 \(\mathfrak{C}\) 是一致的，并且如果形式化算术的公理在 \(\mathfrak{C}\) 中是定理，那么 \(\mathfrak{C}\) 不是范畴性的。他那巧妙方法的基本思想在于，在元数学陈述与形式化算术的某些命题之间建立一一对应关系（当然，是通过“有限程序”）。这里我们给出论证的概要草图（**）。对于 \(\mathfrak{C}\) 中每个作为项或关系的组合体A，我们关联（通过一种可准机械应用的显式构造程序）一个整数 \(g(A)\) ，其方式是一一对应的。同样地，对于 \(\mathfrak{C}\) 中的每一个证明 D（视为一系列装配；参见第一章，§2，第2节），我们可以关联一个整数 \(h(\mathbf{D})\) ，其方式是一一对应的。最后，我们可以给出一个构造关系 \(\mathrm{P}\{x, y, z\}\) （在 \(\mathcal{C} (*)\) 中）的明确步骤，使得在 \(\mathcal{C}, \mathrm{P}\{x, y, z\}\) 蕴含 \(x, y, z\) 是整数，并且满足以下两个条件：

\begin{enumerate}
\def\labelenumi{(\arabic{enumi})}
\tightlist
\item
  如果 \(\mathbf{D}\) 是 \(\mathbf{A}\{\lambda\}\) 的一个证明，其中 \(\mathbf{A}\{x\}\) 是 \(\mathfrak{G}\) 中的一个关系，以及 \(\lambda\) 是一个确定的整数（即 \(\mathfrak{G}\) 中一个作为整数的项），则
\end{enumerate}

\[
\mathbf {P} \{\lambda , g (\mathbf {A} \{x \}), h (\mathbf {D}) \}
\]

在 \(\mathfrak{G}\) 中是定理。

\begin{enumerate}
\def\labelenumi{(\arabic{enumi})}
\setcounter{enumi}{1}
\tightlist
\item
  如果确定的整数 \(\mu\) 不具有 \(h(\mathbf{D})\) 的形式，或者如果 \(\mu = h(\mathbf{D})\) 并且 \(\mathbf{D}\) 不是对 \(\mathbf{A}\{\lambda\}\) 的证明，则（非 \(\mathbf{P}\{\lambda, g(\mathbf{A}\{x\}), \mu\}\) ）在 \(\mathfrak{C}\) 中是定理。
\end{enumerate}

现在令 \(S\{x\}\) 为关系（非 \((\exists z)P\{x, x, z\}\) ），并令 \(\gamma = g(S\{x\})\) ，这是 \(\mathfrak{C}\) 中的一个项。如果 \(\mathfrak{C}\) 并不矛盾，则不存在该命题 \(S\{\gamma\}\) 在 \(\mathfrak{C}\) 中的证明。因为如果D是这样的证明，那么

\[
\mathrm{P} \{\gamma , g (\mathrm{S} \{x \}), h (\mathrm{D}) \}
\]

将是 \(\mathcal{G}\) 中的一个定理。这个关系正是 \(\mathrm{P}\{\gamma, \gamma, h(\mathrm{D})\}\) ，并且因此 \((\exists z)\mathrm{P}\{\gamma, \gamma, z\}\) 也将是 \(\mathcal{G}\) 中的一个定理；但由于最后一个关系等价于（并非 \(\mathrm{S}\{\gamma\}\) ), \(\mathcal{G}\) 将是矛盾的。此外，刚才所述表明，对于每个确定的整数 \(\mu\) ，（非 \(\mathrm{P}\{\gamma, \gamma, \mu\}\) ）在 \(\mathcal{G}\) 中是定理。由此可见，在 \(\mathcal{G}\) 中不存在命题（并非 \(\mathrm{S}\{\gamma\}\) ）的证明，因为该关系等价于 \((\exists z)\mathrm{P}\{\gamma, \gamma, z\}\) ，且存在整数 \(\mu\) 使得 \(\mathrm{P}\{\gamma, \gamma, \mu\}\) 将意味着 \(\mathcal{G}\) 是矛盾的，依据前述论证（ \(\dagger\) ). 哥德尔这一元数学定理后来已在多个方向上得到推广。{[}48{]}，第十一章）（**）。

关系 \(S\{\gamma\}\) 在 \(\mathcal{C}\) 中被证明具有这样的性质：在 \(\mathcal{C}\) 中不存在它或其否定的任何证明；显然，它就是为了适应论证的需要而人为构造的，与任何数学问题都没有自然的联系。更有趣的是，如果 \(\mathcal{C}\) 表示集合论（带有冯·诺依曼-贝尔奈斯公理系统）的理论，连续统假设及其否定在 \(\mathcal{C}\) 中均不可证明。这一非凡结果分两步确立：1940年，哥德尔证明了在 \(\mathcal{C}\) 中加入 \(2^{\aleph_0} = \aleph_1\) 这一假设 {[}并不矛盾{]}；而最近P. Cohen已经证明，当关系 \(2^{\aleph_0} = \aleph_2\) （或 \(2^{\aleph_0} = \aleph_n\) ，对于任意整数 \(n > 1\) ）被附加到 \(\mathcal{C}\) 时，结论也同样成立 {[}49{]}.

判定问题（Entscheidungsproblem）无疑是元数学中最为雄心勃勃的问题。该问题在于，对于给定的形式化语言，是否能够构想出一种准机械的“通用程序”，当将其应用于所考虑的形式系统中的任意关系时，能在有限次运算内指示该关系是否为真。这一问题的解决构成了莱布尼茨宏伟构想的重要组成部分，而且似乎在某一时期，希尔伯特学派曾相信他们已非常接近解决方案。事实上，对于包含少量初始符号和公理的形式系统，确实可以描述出这样的程序（{[}48{]}，第136-141页；参见第一章附录练习7）。但是，通过精确定义“通用程序”的含义来使判定问题精确化的努力，迄今为止只得到了否定的结果（{[}48{]}，第432-439页）。此外，理论G的判定问题的解将立即告诉我们G是否矛盾，因为“通用程序”可以应用于 \(\mathcal{C}\) 中的一个关系及其否定；而且，正如我们将看到的，以这种方式解决该问题的可能性被排除，就通常的数学理论而言（*）。

事实上，正是在数学理论的一致性问题——数学的起源与核心——上，结果被证明是最具欺骗性的。在1920年至1930年间，希尔伯特及其学派发展了解决这些问题的新方法。在证明了某些部分形式系统（涵盖算术的一部分，参见 {[}42{]}, {[}43 c{]}）的一致性之后，他们以为自己已接近胜利，不仅将证明算术的一致性，还将证明集合论的一致性，此时哥德尔利用算术的非范畴性推断出，任何包含算术的理论E，都无法通过希尔伯特的“有限程序”证明其一致性（†）。

然而，哥德尔定理并未完全关闭证明一致性的尝试之门，前提是希尔伯特关于“有限程序”的限制（至少部分地）被放弃。因此，在1936年，根岑 {[}47{]} 成功地通过“直观”地使用超限归纳法直至可数序数 \(\varepsilon_0\) ，证明了形式化算术的一致性（第三章，§6，习题14）(**)。对于此类论证所能赋予的“确定性”程度，无疑低于那些满足希尔伯特原始限制的论证，并且本质上取决于个别数学家的个人心理。但同样正确的是，如果类似的“证明”使用“直觉的”超限归纳法直至某个给定序数，并且能够应用于例如实数理论或集合论的重要部分，那么它们将被视为向前迈出的重要一步。

此外，在集合论本身中，还存在许多与理论中盛行的众多“假设”相关的“相对”一致性问题。该领域最引人注目的成果同样归功于哥德尔。1940年，他证明了，如果以冯·诺依曼-贝尔奈斯系统的公理（除选择公理外）为基础的理论是一致的，那么在这些公理上附加一个非常强的选择公理形式（即使用符号 \(\tau\) 给我们的那种类型）以及广义连续统假设所得到的理论也是一致的 {[}44 b{]}。更近期，I. Novak-Gal 已证明 Zermelo-Fraenkel 系统的一致性蕴含 von Neumann-Bernays-Gödel 系统的一致性（*）。

